\subsection{有限作用窗：从振幅积累算到黄金律误差}
\label{sec:electron-finite-window-worked}

上一节的矩形脉冲给出 $\sinc^2$ 跃迁谱。它来自不同时刻发生的跃迁振幅相加：共振时相位相同，失谐时相位随时间转动。为了比较不同脉冲形状，把这一积分保留下来。设初态到末态通道 $j$ 的相互作用绘景矩阵元为 $\hbar\mathcal G_jw(t)\ee^{\ii\Delta_jt}$，其中 $\mathcal G_j$ 是耦合率，$\Delta_j$ 是失谐，两者单位均为 ${\rm s^{-1}}$；无量纲函数 $w(t)$ 描述耦合何时打开、何时关闭以及怎样变化。由一阶含时微扰论，
\begin{equation}
 A_j^{(1)}=-\ii\mathcal G_j\int\dd t\,w(t)\ee^{\ii\Delta_jt},\qquad
 P_j^{(2)}=|\mathcal G_j|^2\left|\int\dd t\,w(t)\ee^{\ii\Delta_jt}\right|^2.
 \label{eq:electron-window-general-amplitude}
\end{equation}
这里已把玻色占据因子或内部矩阵元吸收到所指定通道的 $\mathcal G_j$ 中。上标表示振幅和概率的微扰阶数，而不是交换光子数。它要求总跃迁概率远小于一；不能在某个单通道概率很小时就忽略所有通道累积造成的耗尽。

矩形窗口 $w(t)=1$（$0<t<T$）给出
\begin{equation*}
 A_j^{(1)}=-\ii\mathcal G_jT\ee^{\ii\Delta_jT/2}
 \sinc(\Delta_jT/2).
\end{equation*}
在共振处各时刻同相累积，振幅正比于 $T$，概率正比于 $T^2$；失谐时较晚的振幅逐渐转向，第一次完全相消出现在 $|\Delta_j|=2\pi/T$。概率峰的角频率全宽半高为 $5.566/T$，相应能量全宽半高为 $5.566\hbar/T$。这个宽度属于有限\emph{相干作用窗}，不是衰变态寿命，也不是谱仪分辨率。

\begin{exbox}[高斯作用窗的线宽与等面积比较]
把矩形开关换成 $w(t)=\exp[-t^2/(2\tau_w^2)]$，求弱跃迁谱，并与具有相同积分面积的矩形窗比较。这里 $\tau_w$ 是振幅包络的时间尺度。

对指数配方，或直接使用高斯 Fourier 积分，有
\begin{equation*}
 \int_{-\infty}^{\infty}\dd t\,
 \ee^{-t^2/(2\tau_w^2)+\ii\Delta t}
 =\sqrt{2\pi}\tau_w\ee^{-\Delta^2\tau_w^2/2}.
\end{equation*}
因此
\begin{equation*}
 A_j^{(1)}=-\ii\mathcal G_j\sqrt{2\pi}\tau_w
 \ee^{-\Delta_j^2\tau_w^2/2},\qquad
 P_j^{(2)}=2\pi|\mathcal G_j|^2\tau_w^2
 \ee^{-\Delta_j^2\tau_w^2}.
\end{equation*}
令谱值降到共振峰的一半，取对数得 $|\Delta_j|\tau_w=\sqrt{\ln2}$，故角频率全宽半高为 $2\sqrt{\ln2}/\tau_w$。它来自\emph{概率}包络，不能直接使用振幅包络的半高宽。

若矩形窗长度取 $T=\sqrt{2\pi}\tau_w$，两窗具有相同积分面积和相同共振概率 $|\mathcal G_j|^2T^2$。高斯窗的谱全宽为 $2\sqrt{2\pi\ln2}/T$；矩形窗则为 $4x_{1/2}/T$，其中 $x_{1/2}$ 是 $\sinc^2x=1/2$ 的首个正根。两者都按 $T^{-1}$ 缩窄，但高斯谱没有旁瓣，矩形谱有代数衰减的旁瓣。延长窗口时，弱跃迁条件仍须独立检查，不能只比较线宽。
\end{exbox}

\begin{exbox}[双脉冲的 Ramsey 条纹：延迟如何进入分辨率]
考虑两段等长矩形作用窗，长度为 $\tau_p$，中心分别在 $-D_p/2$ 与 $D_p/2$，其中 $D_p\ge\tau_p$。第二段耦合比第一段多一个可控相位 $\phi_p$，即使用复开关
\begin{equation*}
 w(t)=\mathbf1_{|t+D_p/2|<\tau_p/2}
       +\ee^{\ii\phi_p}\mathbf1_{|t-D_p/2|<\tau_p/2}.
\end{equation*}
符号 $\mathbf1$ 在指定区间内取一、区间外取零。假定两脉冲之间相干保持，并且总跃迁仍可用一阶振幅处理。

每段积分都是同一个 $\tau_p\sinc(\Delta\tau_p/2)$，只差中心时刻与参考相位。先相加振幅，得
\begin{align*}
 A^{(1)}&=-\ii\mathcal G\tau_p\sinc(\Delta\tau_p/2)
 [\ee^{-\ii\Delta D_p/2}+\ee^{\ii\phi_p+\ii\Delta D_p/2}]\\
 &=-2\ii\mathcal G\tau_p\sinc(\Delta\tau_p/2)
 \ee^{\ii\phi_p/2}\cos[(\Delta D_p+\phi_p)/2].
\end{align*}
因此
\begin{equation*}
 P^{(2)}=4|\mathcal G|^2\tau_p^2\sinc^2(\Delta\tau_p/2)
 \cos^2[(\Delta D_p+\phi_p)/2].
\end{equation*}
单脉冲时宽决定宽包络，脉冲中心间隔决定包络内的细条纹。$\phi_p=0$ 且 $D_p\gg\tau_p$ 时，中央条纹内的 $\sinc$ 近似为一；由 $\cos^2(\Delta D_p/2)=1/2$ 得全宽半高约为 $\pi/D_p$。因此长延迟可以缩窄条纹，而不增加两段实际作用时间。

若 $\phi_p$ 在重复实验间均匀随机，$\overline{\cos^2}=1/2$，平均概率退化为两个单脉冲概率之和，细条纹消失。把概率先相加相当于预先作了这一失相干假设；它不是相干双脉冲实验的一般规则。
\end{exbox}

单个共振末态的概率随 $T^2$ 增长，而黄金律给出的总概率随 $T$ 增长。要看清这一变化，需要把连续末态在不同失谐处的贡献积分起来。将末态数目与耦合权重合并，定义所选初态的\emph{跃迁谱密度}
\begin{equation}
 \mathcal J_{\rm tr}(\Delta)=\sum_j|\mathcal G_j|^2\delta(\Delta-\Delta_j),\qquad
 P^{(2)}(T)=\int\dd\Delta\,\mathcal J_{\rm tr}(\Delta)
 T^2\sinc^2(\Delta T/2).
 \label{eq:electron-window-transition-spectrum}
\end{equation}
$\mathcal J_{\rm tr}$ 的单位为 $\mathrm{s^{-1}}$，它同时包含末态数目与每个末态的耦合强度。窗口变长时，共振峰高度按 $T^2$ 增长，但参与的连续末态频率宽度按 $1/T$ 缩小，两者相乘才给线性的总概率。对孤立离散态没有这一连续谱积分，不能直接使用恒定黄金律速率。

为了算出有限时间误差，取以共振为中心的连续谱模型
\begin{equation*}
 \mathcal J_{\rm tr}(\Delta)=\frac{J_0}{1+\Delta^2\tau_c^2},
 \qquad J_0>0,\qquad \tau_c>0.
\end{equation*}
$J_0$ 是谱峰值，$\tau_c$ 是该谱宽倒数所定义的相关时间。失谐轴延拓为整个实轴，适用于实际阈值远离所采样共振窗的教学模型。积分可完整算成
\begin{equation}
 P^{(2)}(T)=\Gamma_0[T-\tau_c(1-\ee^{-T/\tau_c})],\qquad
 \Gamma_0=2\pi J_0.
 \label{eq:electron-window-lorentzian-probability}
\end{equation}
短时间给 $P^{(2)}\simeq\Gamma_0T^2/(2\tau_c)$；长时间给 $P^{(2)}\simeq\Gamma_0(T-\tau_c)$。黄金律 $P\simeq\Gamma_0T$ 因而同时需要 $\tau_c/T\ll1$ 和 $\Gamma_0T\ll1$：前者使连续谱有时间去相位，后者使初态尚未明显耗尽。

\begin{exbox}[黄金律的适用时间窗是否存在]
对上述 Lorentz 跃迁谱，给定允许的有限时间相对误差 $\eta$ 和最大耗尽概率 $p_*$，求一组便于手算的充分条件。

先无量纲化，令 $x=T/\tau_c$、$\epsilon_c=\Gamma_0\tau_c$。精确二阶结果与黄金律之比为
\begin{equation*}
 R(x)=\frac{P^{(2)}(T)}{\Gamma_0T}
 =1-\frac{1-\ee^{-x}}{x},\qquad
 1-R(x)=\frac{1-\ee^{-x}}{x}<\frac1x.
\end{equation*}
这里把相对误差定义为相对于黄金律预测的偏差。求导得
\begin{equation*}
 R'(x)=\frac{1-(1+x)\ee^{-x}}{x^2}>0,
\end{equation*}
因为 $\ee^x>1+x$。因此增加作用时间会单调改善黄金律的有限窗误差，却同时增加耗尽。

由上界，$T\ge\tau_c/\eta$ 足以使相对误差不超过 $\eta$；又因 $P^{(2)}<\Gamma_0T$，取 $T\le p_*/\Gamma_0$ 足以控制耗尽。两条件能够同时满足，只需
\begin{equation*}
 \frac{\tau_c}{\eta}\le T\le\frac{p_*}{\Gamma_0},\qquad
 \Gamma_0\tau_c\le\eta p_*.
\end{equation*}
这是一组充分条件，未必是最宽的时间窗。它直接说明：只有环境去相位时间显著短于初态耗尽时间，才存在可用的恒定速率区。若区间为空，应回到有限时间概率或完整动力学，而不是强行定义一个黄金律速率。
\end{exbox}

\begin{derivation}[直接积分连续谱，而不预先代入能量守恒]
写 $T^2\sinc^2(\Delta T/2)=\int_0^T\dd t\int_0^T\dd t'\,\ee^{\ii\Delta(t-t')}$，并使用 Lorentz 线形的 Fourier 积分
\begin{equation*}
 \int_{-\infty}^{\infty}\dd\Delta\,
 \frac{J_0\ee^{\ii\Delta u}}{1+\Delta^2\tau_c^2}
 =\frac{\pi J_0}{\tau_c}\ee^{-|u|/\tau_c}.
\end{equation*}
两时间正方形区域中，相同时间差 $u>0$ 对应的线段长度为 $T-u$，故
\begin{align*}
 P^{(2)}(T)&=\frac{2\pi J_0}{\tau_c}
 \int_0^T(T-u)\ee^{-u/\tau_c}\dd u\\
 &=2\pi J_0[T-\tau_c(1-\ee^{-T/\tau_c})].
\end{align*}
因此恒定速率是有限时间积分的渐近结果。谱阈值、窄峰或离散末态改变了被积分的 $\mathcal J_{\rm tr}$ 时，应先重算这个积分，再判断是否存在黄金律区。
\end{derivation}
